\begin{pastexam}[label={px:gw3-110-ds-05}]{關務三等資訊處理 110 年 第五題}
（申論；30 分）

\begin{itemize}[leftmargin=2.2em]
\item[㈠] 如下圖設背包限重 100，有 A、B、C、D、E 共五個不可分割物件，請問依貪婪策略（Greedy Algorithm），0\_1 整數
背包問題（knapsack problem）/貨物裝載問題（cargo loading problem）其最大利益為何？其對應的 0\_1 整數規劃為何？（20 分）
\item[㈡] 有 A、B、C、D、E 共五個可分割物件，請問依貪婪策略，0\_1 分數背包其最大利益為何？（10 分）
\end{itemize}
\begin{center}
\begin{tabular}{|c|c|c|}
\hline
物件 & 重量 & 利益 \\
\hline
A & 10 & 20 \\
\hline
B & 20 & 30 \\
\hline
C & 30 & 66 \\
\hline
D & 40 & 40 \\
\hline
E & 50 & 60 \\
\hline
\end{tabular}
\end{center}
\end{pastexam}

\begin{pastsoln}
答案：㈠ 最大利益 $156$（取 A、B、C、D）；整數規劃見下。 ㈡ 最大利益 $164$（A、B、C 全取，E 取 $4/5$）。

㈡的「0\_1 分數背包」指的是物件可以分割的分數背包（fractional knapsack），以下照這個意思解。 兩小題的貪婪策略
都依「單位重量的利益」由大到小挑選（資結書第~19.2~節）：
\begin{center}
\begin{tabular}{lccccc}
\toprule
物件 & A & B & C & D & E \\
\midrule
重量 $w_i$ & 10 & 20 & 30 & 40 & 50 \\
利益 $p_i$ & 20 & 30 & 66 & 40 & 60 \\
$p_i / w_i$ & 2.0 & 1.5 & 2.2 & 1.0 & 1.2 \\
排名 & 2 & 3 & 1 & 5 & 4 \\
\bottomrule
\end{tabular}
\end{center}

\textbf{㈠ 不可分割（0/1 背包）。} 依 C、A、B、E、D 的順序，放得下就放：
\begin{itemize}
\item C：重 $30 \le 100$，放入。 已用 $30$，利益 $66$。
\item A：重 $10$，放入。 已用 $40$，利益 $86$。
\item B：重 $20$，放入。 已用 $60$，利益 $116$。
\item E：重 $50$，$60 + 50 = 110 > 100$，放不下，跳過。
\item D：重 $40$，$60 + 40 = 100$，剛好放入。 利益 $156$。
\end{itemize}
貪婪的結果是 A、B、C、D，總重 $100$，利益 $156$。

\textbf{它是不是最佳？} 貪婪法對 0/1 背包不保證最佳（資結書第~19.5~節），要另外檢查。 五件全拿總重 $150$、總利益
$216$，超重 $50$，所以至少要拿掉總重 $\ge 50$ 的物件；問題變成「拿掉的物件總重 $\ge 50$ 時，損失的利益最少是多少」。
拿掉 E（重 $50$，損失 $60$）、或拿掉 A 和 D（重 $50$，損失 $60$），損失都是 $60$；其他總重 $\ge 50$ 的組合，例如
B 和 D（損失 $70$）、B 和 C（損失 $96$）、A、B 和 D（損失 $90$）、含 C 的組合（C 本身就損失 $66$）都更多，而總重
$\ge 50$ 的單一物件只有 E。 所以最少損失 $60$，最大利益 $216 - 60 = 156$，貪婪的結果正好是最佳解之一（另一個最佳解
是 B、C、E）。 也可以用資結書第~18.5~節的動態規劃驗證，同樣得到 $156$。

\textbf{對應的 0/1 整數規劃。} 令 $x_A, \ldots, x_E \in \{0, 1\}$，$x_i = 1$ 表示放入物件 $i$：
\begin{align*}
\text{maximize}\quad & 20x_A + 30x_B + 66x_C + 40x_D + 60x_E \\
\text{subject to}\quad & 10x_A + 20x_B + 30x_C + 40x_D + 50x_E \le 100, \\
& x_A, x_B, x_C, x_D, x_E \in \{0, 1\}.
\end{align*}
貪婪的解是 $(x_A, x_B, x_C, x_D, x_E) = (1, 1, 1, 1, 0)$，目標值 $156$。

\textbf{㈡ 可分割（分數背包）。} 同樣依 C、A、B、E、D 的順序，放得下就整件放，放不下就放入剩餘容量那麼多：
\begin{itemize}
\item C、A、B 整件放入：已用 $60$，利益 $116$，剩 $40$。
\item E 重 $50 > 40$，放入 $40/50 = 4/5$ 件，利益 $60 \times 4/5 = 48$。 背包滿了。
\item D 放不進去。
\end{itemize}
最大利益 $116 + 48 = 164$，對應 $(x_A, x_B, x_C, x_D, x_E) = (1, 1, 1, 0, 0.8)$。 分數背包依單位利益挑選一定是最佳解
（資結書定理~19.2），所以 $164$ 就是最大值。 它也是㈠的上界：0/1 背包的解也是分數背包的可行解，$156 \le 164$。
\end{pastsoln}
